\section{Entender el lenguaje a partir de bytes: costos y beneficios de eliminar el tokenizer}

La sección anterior mostró que las output queries pueden producir secuencias; esta sección aplica la interfaz unificada al lenguaje. Un Transformer habitual usa primero un tokenizer para dividir la cadena en subwords; esto acorta mucho la secuencia, pero incorpora un vocabulario, escrituras y reglas específicas de cada idioma. Gracias al latent bottleneck, Perceiver IO puede procesar una raw byte sequence más larga e intenta debilitar los supuestos del tokenizer.

\lecturefigure{slide-22-language-tokenization.jpg}{Lenguaje, Transformer y tokenización: una misma oración puede dividirse en caracteres, bytes o subwords.}{Video oficial actual de Stanford Online, 00:43:44--00:45:13.}

\begin{knowledgebox}{Qué es el byte-level input}
El texto se codifica primero como bytes UTF-8; cada valor de byte está entre 0 y 255. Es un alfabeto fijo compartido por todos los idiomas y escrituras, y evita las palabras desconocidas; el costo es que un carácter puede ocupar varios bytes, y la secuencia resulta mucho más larga que con la tokenización en subwords.
\end{knowledgebox}

\subsection{Por qué eliminar el tokenizer}

Esta sección examina el tokenizer como un sesgo lingüístico que puede elegirse. El vocabulario del tokenizer funciona bien en idiomas con muchos recursos, pero puede fragmentar en exceso idiomas poco frecuentes, ortografías, código, emoji o palabras nuevas. Distintos idiomas pueden requerir distintos vocabulary, y un sistema multilingüe además debe manejar diferencias de normalización y de presegmentación en palabras. Los bytes ofrecen el nivel más bajo de símbolos unificados, de modo que el modelo aprende por sí mismo las fronteras.

\lecturefigure{slide-23-why-remove-tokenizers.jpg}{Dependencia del idioma del tokenizer y representación en bytes UTF-8.}{Video oficial actual de Stanford Online, 00:45:14--00:45:37.}

\teachervoice{El docente no dijo que el tokenizer sea «erróneo»; enfatizó que es un inductive bias añadido de forma manual. Los bytes son más generales, pero trasladan al modelo el aprendizaje de la estructura y el costo de las secuencias largas; si vale la pena depende de la escala de datos, la cobertura de idiomas y el presupuesto de cómputo.}

\begin{warningbox}{Token-free no significa cero preprocesamiento}
La Unicode normalization, la codificación UTF-8, el padding, el masking y la byte position siguen siendo decisiones de preprocesamiento. Un byte vocabulary fijo solo elimina la subword segmentation; no elimina la normalización del texto ni la limpieza de datos.
\end{warningbox}

\subsection{Masked Language Modeling y fine-tuning en GLUE}

El entrenamiento sigue el masked language modeling al estilo de BERT: se enmascara una parte de los bytes de entrada y las decoder queries predicen el objetivo en las posiciones correspondientes; después se hace fine-tuning en las tareas de GLUE añadiendo sentence/task queries. Para el conjunto enmascarado $\mathcal{M}$, el objetivo puede escribirse como:

\[
\mathcal{L}_{\text{MLM}}=-\sum_{t\in\mathcal{M}}\log p_\theta(b_t\mid b_{\setminus\mathcal{M}}).
\]

Aquí, $b_t$ es el byte real en la posición $t$; $b_{\setminus\mathcal{M}}$ es el contexto no enmascarado; $p_\theta$ es la distribución de salida del modelo; $\mathcal{M}$ son las masked positions.

\lecturefigure{slide-24-masked-language-modeling.jpg}{Byte-level masked language modeling de Perceiver IO.}{Video oficial actual de Stanford Online, 00:45:38--00:45:53.}

Durante el fine-tuning, la tarea puede decodificar logits de sentence, entailment, grammar o paraphrase mediante queries adicionales. El encoder/latent processor compartido recibe bytes, y la interface de salida cambia según la tarea.

\lecturefigure{slide-25-glue-finetuning.jpg}{Del byte-level MLM a GLUE: distintas output queries decodifican los logits de cada tarea.}{Video oficial actual de Stanford Online, 00:45:54--00:47:17.}

\begin{importantbox}{Qué aporta el latent bottleneck a los bytes}
La byte sequence es muy larga, pero la deep self-attention se realiza solo sobre latents fijos; la cross-attention de entrada es lineal en el número de bytes. El modelo cambia un cuello de botella de cómputo por un vocabulary unificado e interioriza el aprendizaje de las fronteras de segmentación.
\end{importantbox}

\subsection{Resultados FLOP-matched y attention diagnostics}

La tabla compara BERT, CANINE y Perceiver IO con distintas profundidades y números de parámetros. Con FLOPs similares, el Perceiver IO byte-level alcanza puntajes de GLUE competitivos; aumentar el tamaño del modelo sigue mejorándolos. Una lectura correcta debe considerar a la vez el tokenizer, la profundidad, los parámetros, los FLOPs y los datos de entrenamiento, en lugar de elegir un solo puntaje.

\lecturefigure{slide-26-language-from-bytes-results.jpg}{Entender el lenguaje directamente a partir de bytes: comparación de parámetros, FLOPs y GLUE entre modelos tokenizados y token-free.}{Video oficial actual de Stanford Online, 00:47:18--00:47:59.}

\begin{knowledgebox}{Lectura de la tabla: igualar FLOPs no equivale a igualar latencia}
Una byte sequence larga cambia el memory access, el kernel de cross-attention y el tamaño de lote. Los mismos FLOPs teóricos pueden dar un wall-clock distinto en distinto hardware; también deben compararse la longitud de secuencia, el rendimiento, la memoria pico y la cobertura de idiomas.
\end{knowledgebox}

La visualización de attention muestra que el modelo atiende a la puntuación, las mayúsculas y minúsculas, los fragmentos de palabras y los byte patterns repetidos. Esto indica que el modelo puede descubrir estructura a partir de símbolos de bajo nivel, pero el attention weight no es una explicación causal ni demuestra que una head «entienda la gramática».

\lecturefigure{slide-27-language-attention-patterns.jpg}{Diagnóstico de los attention patterns de lenguaje en el Perceiver IO byte-level.}{Video oficial actual de Stanford Online, 00:48:00--00:48:43.}

\teachervoice{El ponente presenta estas figuras como una comprobación de que «el modelo efectivamente atiende a estructuras reconocibles», no como una conclusión completa de interpretabilidad. Sobre un attention map deben hacerse también pruebas de ablation, counterfactual y sensibilidad de la salida.}

\begin{warningbox}{La attention visualization se presta a narrativas exageradas}
Los patrones de color pueden sugerir hipótesis, pero por sí solos no explican cómo influye la información en la predicción final. Varias heads, la residual y el MLP actúan en conjunto; un pattern atractivo puede tener una contribución muy débil a la tarea.
\end{warningbox}

\subsection{Resumen de la sección}

El Perceiver IO byte-level cambia secuencias largas y más aprendizaje de estructura por un vocabulary unificado entre idiomas. Los resultados de MLM/GLUE demuestran que el enfoque es viable, pero no niegan de forma general las ventajas del tokenizer en eficiencia de datos y despliegue, y el attention map es solo evidencia de diagnóstico.
